\documentclass[11pt]{article}
\usepackage[margin=0.88in]{geometry}
\usepackage{amsmath,amssymb,lmodern}
\usepackage[T1]{fontenc}
\usepackage{needspace}
\usepackage[colorlinks=true,linkcolor=blue,urlcolor=blue]{hyperref}
\setlength{\emergencystretch}{3em}
\setlength{\parskip}{0.35em}
\setcounter{secnumdepth}{0}
\providecommand{\tightlist}{\setlength{\itemsep}{0pt}\setlength{\parskip}{0pt}}
\title{Torsion transplantation and a uniform reverse Kohler--Jobin inequality for Robin problems}
\author{Henry Zweiman}
\date{September 15, 2026}
\begin{document}
\maketitle
\section{Abstract}\label{abstract}

We prove that balls maximize the product of the first Robin eigenvalue
and a sufficiently large power of the torsional rigidity among bounded
planar Lipschitz sets of prescribed area. The exponent threshold is
independent of the positive Robin parameter and the area. This answers
the planar case of the reverse Kohler-Jobin question posed by Buttazzo,
Cito and Solombrino. More generally, the argument applies to the Robin
\(p\)-Laplacian in dimension \(n\) when \(1<p\leq n/(n-1)\), using the
integral of the \(p\)-torsion function. The main step is a composition
identity with a nonpositive boundary correction, which transplants the
ball eigenfunction through the torsion function and preserves the
reference energy. The known pointwise Talenti comparison then yields a
linear upper bound for the spectral deficit in terms of the torsional
deficit. An elementary global bound completes the optimization argument.
For larger exponents, the product deficit is equivalent to the torsional
deficit, uniformly in the Robin parameter. The higher-dimensional linear
Robin problem is not resolved.

\section{1. Introduction and main
results}\label{introduction-and-main-results}

For a fixed positive Robin parameter, the ball minimizes the first
eigenvalue and maximizes torsional rigidity among sets of prescribed
volume. These two inequalities act in opposite directions on their
product. Increasing the power of torsion suggests that its maximizer
should eventually maximize the product, but the separate isoperimetric
inequalities do not prove this assertion: one needs quantitative control
of the competing increase in the eigenvalue.

Buttazzo, Cito and Solombrino {[}BCS, Section 5, Q2{]} ask whether the
ball maximizes the first Robin eigenvalue times a sufficiently large
power of Robin torsional rigidity. The corresponding Dirichlet reverse
Kohler-Jobin inequality is established by Bucur, Lamboley, Nahon and
Prunier {[}BLNP, Corollary 1.6{]}. Here we answer the Robin question in
the plane. Our exponent threshold is uniform in both the positive
boundary parameter and the prescribed area. The proof also gives a
nonlinear extension in the range of the known pointwise Robin
rearrangement theorem.

Let \(n\geq2\), \(p>1\) and \(\beta>0\). For a nonempty bounded open set
\(\Omega\subset\mathbb R^n\) with Lipschitz boundary, define

\[\begin{aligned}
\mathcal E_{\beta,p,\Omega}(f)&=\int_\Omega|\nabla f|^p\,dx+\beta\int_{\partial\Omega}|f|^p\,d\mathcal H^{n-1},\\
\lambda_{\beta,p}(\Omega)&=\inf_{f\in W^{1,p}(\Omega)\setminus\{0\}}
\frac{\mathcal E_{\beta,p,\Omega}(f)}{\int_\Omega|f|^p\,dx}.
\end{aligned}\tag{1.1}\]

Let \(w=w_{\beta,p,\Omega}\) be the nonnegative weak solution of

\[-\Delta_p w=1\quad\text{in }\Omega,
\qquad |\nabla w|^{p-2}\partial_\nu w+\beta w^{p-1}=0
\quad\text{on }\partial\Omega,\tag{1.2}\]

where \(\Delta_p w=\operatorname{div}(|\nabla w|^{p-2}\nabla w)\), and
put

\[\tau_{\beta,p}(\Omega)=\int_\Omega w\,dx.\tag{1.3}\]

Our convention for \(\tau\) is the integral of \(w\), rather than its
\((p-1)\)-st power. For \(p=2\) it is the usual torsional rigidity
\(T_\beta\). Let \(B\) be a ball with \(|B|=|\Omega|=V\). We omit
subscripts when the parameters are fixed and write

\[\theta(\Omega)=\frac{\tau(\Omega)}{\tau(B)},\qquad
\mathcal R_q(\Omega)=\frac{\lambda(\Omega)\tau(\Omega)^q}
{\lambda(B)\tau(B)^q}.\tag{1.4}\]

\Needspace{6\baselineskip}\textbf{Theorem 1.1 (uniform reverse Robin inequality).} Fix \(n\geq2\)
and \(1<p\leq n/(n-1)\). There is a finite number
\(Q(n,p)>\max\{1,p-1\}\) such that, for every \(\beta>0\), \(V>0\) and
nonempty bounded open Lipschitz set \(\Omega\) of volume \(V\),

\[\lambda_{\beta,p}(\Omega)\tau_{\beta,p}(\Omega)^q
\leq\lambda_{\beta,p}(B)\tau_{\beta,p}(B)^q
\qquad(q\geq Q(n,p)).\tag{1.5}\]

For every such \(q\),

\[1-\theta(\Omega)\leq1-\mathcal R_q(\Omega)
\leq q\bigl(1-\theta(\Omega)\bigr).\tag{1.6}\]

Equality in (1.5) holds if and only if \(\Omega\) is a translate of
\(B\).

In particular, \(Q(2,2)\) supplies a single exponent threshold for the
planar problem {[}BCS, Q2{]}, simultaneously for all positive Robin
parameters. Neither convexity nor connectedness is required. Formula
(1.6) says that, for a fixed admissible exponent, the normalized product
deficit and torsional deficit are equivalent with constants independent
of \(\beta\), \(V\) and \(\Omega\).

The estimate underlying the theorem is the following one-sided spectral
control.

\Needspace{6\baselineskip}\textbf{Theorem 1.2 (linear spectral control).} Under the hypotheses of
Theorem 1.1, there is \(C(n,p)<\infty\), independent of \(\beta\) and
\(V\), such that

\[0\leq\frac{\lambda(\Omega)}{\lambda(B)}-1
\leq C(n,p)\bigl(1-\theta(\Omega)\bigr)
\quad\text{if}\quad
1-\theta(\Omega)\leq C(n,p)^{-1}.\tag{1.7}\]

The essential rearrangement input is {[}AGM, Theorem 1.2(i){]}:

\[w_\Omega^\#\leq w_B\quad\text{a.e. on }B
\qquad\left(1<p\leq\frac n{n-1}\right),\tag{1.8}\]

where \(\#\) denotes Schwarz decreasing rearrangement. For \(n=p=2\),
this is the comparison of Alvino, Nitsch and Trombetti {[}ANT{]}. We use
their theorem with constant forcing and exactly the same Robin parameter
on both sets. The only dimensional restriction in our proof occurs at
(1.8).

We also use the Robin \(p\)-Faber-Krahn inequality and rigidity {[}BD,
Theorem 1.1{]}. Its extension from connected domains to the present
class is elementary. A bounded Lipschitz set has finitely many
components \(\Omega_j\), each Lipschitz, and its first eigenvalue is
\(\min_j\lambda(\Omega_j)\). The eigenvalue of a ball strictly decreases
with its radius at fixed \(\beta\): dilation by \(t>1\) changes a trial
quotient with gradient part \(A\) and boundary part \(D\) into
\(t^{-p}A+t^{-1}D\), which is at most \(t^{-1}(A+D)\). Consequently
\(\lambda(tB)\leq t^{-1}\lambda(B)<\lambda(B)\). Apply {[}BD{]} to every
component and then enlarge its comparison ball to volume \(V\). This
proves \(\lambda(\Omega)\geq\lambda(B)\), with strict inequality if
\(\Omega\) is disconnected; the connected equality case is precisely
{[}BD{]}. The positivity and simplicity of the ball ground state are
standard parts of this Robin eigenvalue theory.

\section{2. A composition identity with a favorable boundary
term}\label{a-composition-identity-with-a-favorable-boundary-term}

\Needspace{6\baselineskip}\textbf{Lemma 2.1.} Suppose \(w\) is a bounded nonnegative weak Robin
\(p\)-torsion function on a bounded Lipschitz set, with
\(0\leq w\leq M\). Let \(h:[0,M]\to[0,\infty)\) be \(C^1\), with
\(h(0)=0\), and set

\[H(s)=\int_0^s|h'(t)|^p\,dt,
\qquad G(s)=h(s)^p-s^{p-1}H(s).\tag{2.1}\]

Then \(h(w)\in W^{1,p}(\Omega)\) and

\[\mathcal E_{\beta,p,\Omega}(h(w))
=\int_\Omega H(w)\,dx+\beta\int_{\partial\Omega}G(w)\,d\mathcal H^{n-1}
\leq\int_\Omega H(w)\,dx.\tag{2.2}\]

\textbf{Proof.} Extend \(h\) and \(H\) as Lipschitz functions beyond
their compact intervals. Sobolev composition and the trace identity
justify testing (1.2) with \(H(w)\). Since \(H'=|h'|^p\), this gives

\[\int_\Omega|\nabla h(w)|^p\,dx
=\int_\Omega H(w)\,dx-\beta\int_{\partial\Omega}w^{p-1}H(w)\,d\mathcal H^{n-1}.\]

Adding the Robin energy of the trace gives the identity.
Hölder\textquotesingle s inequality yields

\[h(s)^p=\left|\int_0^s h'(t)\,dt\right|^p
\leq s^{p-1}\int_0^s|h'(t)|^p\,dt,\]

so \(G\leq0\). The interior bounds on \(w\) hold for its Sobolev trace
as well: apply the trace operator to the truncations \((w-M)_+\) and
\(w_-\). Thus all boundary terms above are integrable. This proves
(2.2). \(\square\)

If \(h\) is linear on \([0,b]\), then \(G(b)=0\). We will impose
precisely this extension below the boundary value of the ball torsion.
It makes the boundary inequality an equality on the reference ball,
without requiring the trace of the competing torsion function to be
constant.

\section{3. Encoding the ball
eigenfunction}\label{encoding-the-ball-eigenfunction}

Let \(B=B_R\) and write \(p'=p/(p-1)\). The radial torsion function
\(v=w_B\) is

\[v(r)=b+\frac{p-1}{p}n^{-1/(p-1)}(R^{p'}-r^{p'}),
\qquad b=\left(\frac R{n\beta}\right)^{1/(p-1)}.\tag{3.1}\]

Let \(\phi\) be the positive first eigenfunction on \(B\), normalized by
\(\int_B\phi^p\,dx=1\), and write \(\lambda_B=\lambda(B)\). Simplicity
and rotational invariance make \(\phi\) radial. Integrating its radial
equation gives

\[-\phi'(r)=\left[\lambda_B r^{1-n}\int_0^r s^{n-1}\phi(s)^{p-1}\,ds\right]^{1/(p-1)}.\tag{3.2}\]

In particular, \(\phi\) strictly decreases for \(r>0\). Put \(M=v(0)\)
and define

\[h(v(r))=\phi(r)\quad(0\leq r\leq R),
\qquad h(s)=\frac{\phi(R)}b s\quad(0\leq s\leq b).\tag{3.3}\]

For \(0<r\leq R\),

\[h'(v(r))=
\left[n\lambda_B r^{-n}\int_0^r s^{n-1}\phi(s)^{p-1}\,ds\right]^{1/(p-1)}.\tag{3.4}\]

The derivative has the finite positive limit
\(\lambda_B^{1/(p-1)}\phi(0)\) at the center. At \(r=R\), the two Robin
boundary conditions give \(h'(b)=\phi(R)/b\), so the pieces in (3.3)
join in \(C^1\). The limiting derivative also verifies differentiability
at \(M\), for example by the mean value theorem. Thus \(h\) is
increasing and \(C^1\) on \([0,M]\).

Its derivative is nondecreasing in its argument. Indeed, the radial
average in (3.4) decreases with \(r\), since \(\phi\) decreases, while
\(v(r)\) decreases. Therefore both \(h\) and \(h'\) take their maxima at
\(M\), and

\[L:=\max_{0\leq s\leq M}(h(s)^p)'
=p\lambda_B^{1/(p-1)}\phi(0)^p.\tag{3.5}\]

For the function \(H\) in Lemma 2.1, the linear extension makes
\(G(b)=0\). Applying the identity on \(B\) gives

\[\int_B H(v)\,dx=\mathcal E_{\beta,p,B}(\phi)=\lambda_B,
\qquad \int_B h(v)^p\,dx=1.\tag{3.6}\]

Only first derivatives of \(h\) are used, including at the center of the
ball.

\section{4. From torsional deficit to global product
optimality}\label{from-torsional-deficit-to-global-product-optimality}

This section assumes (1.8). Define

\[\delta=\tau(B)-\tau(\Omega),\qquad c=L\tau(B),
\qquad\theta=\frac{\tau(\Omega)}{\tau(B)}.\tag{4.1}\]

The comparison ensures \(0<\theta\leq1\) and bounds \(w\) by \(M\).
Lemma 2.1, equimeasurability and the monotonicity of \(H\) imply

\[\mathcal E_{\beta,p,\Omega}(h(w))
\leq\int_\Omega H(w)\,dx
=\int_B H(w^\#)\,dx
\leq\int_B H(v)\,dx=\lambda_B.\tag{4.2}\]

The denominator satisfies

\[\begin{aligned}
0\leq1-\int_\Omega h(w)^p\,dx
&=\int_B\bigl(h(v)^p-h(w^\#)^p\bigr)\,dx\\
&\leq L\int_B(v-w^\#)\,dx=L\delta.
\end{aligned}\tag{4.3}\]

The trial function is nonzero, since \(h\) is strictly increasing and
\(w\) is nonzero and nonnegative. Its Rayleigh quotient yields

\[\frac{\lambda(\Omega)}{\lambda_B}
\leq\frac1{1-c(1-\theta)}
\quad\text{if }c(1-\theta)<1.\tag{4.4}\]

In particular, Faber-Krahn and (4.4) imply

\[0\leq\frac{\lambda(\Omega)}{\lambda_B}-1
\leq2c(1-\theta)
\quad\text{if }1-\theta\leq\frac1{2c}.\tag{4.5}\]

Since \(h(0)=0\) and \(h^p\) has Lipschitz constant \(L\), (3.6) gives
\(1\leq L\tau(B)\), hence \(c\geq1\).

We also need a global bound. Testing the torsion equation with \(w\)
gives \(\mathcal E(w)=\tau(\Omega)\). Using \(w\) in the Rayleigh
quotient and then applying Hölder gives

\[\lambda(\Omega)\leq\frac{\tau(\Omega)}{\int_\Omega w^p\,dx}
\leq\frac{V^{p-1}}{\tau(\Omega)^{p-1}}.\tag{4.6}\]

This estimate alone does not identify the maximizing ball, but it
controls domains outside the neighborhood covered by (4.5).

\Needspace{6\baselineskip}\textbf{Proposition 4.1 (a threshold from ball data).} Set

\[\begin{aligned}
k&=\frac{\lambda_B\tau(B)^{p-1}}{V^{p-1}},\qquad t_0=1-\frac1{2c},\\
q_0&=\max\left\{2c,\ p-1+\frac{\log(1/k)}{-\log t_0}\right\}.
\end{aligned}\tag{4.7}\]

Then \(\mathcal R_q(\Omega)\leq1\) for all \(q\geq q_0\). For all
\(q\geq q_0+1\),

\[1-\theta\leq1-\mathcal R_q(\Omega)\leq q(1-\theta).\tag{4.8}\]

\textbf{Proof.} By (4.6), \(0<k\leq1\), and \(1/2\leq t_0<1\). If
\(\theta\geq t_0\), then (4.4) gives

\[\log\frac{\lambda(\Omega)}{\lambda_B}
\leq-\log(1-c(1-\theta))
\leq2c(1-\theta)\leq-2c\log\theta.\]

Thus \(\lambda(\Omega)/\lambda_B\leq\theta^{-2c}\), proving the product
bound for \(q\geq q_0\). If \(\theta\leq t_0\), (4.6) gives instead

\[\mathcal R_q(\Omega)\leq k^{-1}\theta^{q-(p-1)}
\leq k^{-1}t_0^{q-(p-1)}\leq1\qquad(q\geq q_0),\]

by the definition of \(q_0\). This covers every \(\theta\in(0,1]\).

For \(q\geq q_0+1\), apply the product bound at \(q_0\) and use
Faber-Krahn to obtain

\[\theta^q\leq\mathcal R_q(\Omega)
=\mathcal R_{q_0}(\Omega)\theta^{q-q_0}
\leq\theta^{q-q_0}\leq\theta.\]

Bernoulli\textquotesingle s inequality \(1-\theta^q\leq q(1-\theta)\)
proves (4.8). \(\square\)

\section{5. Uniformity in the Robin
parameter}\label{uniformity-in-the-robin-parameter}

The constants in Proposition 4.1 involve only the reference ball. The
next step shows that their dependence on its radius and on \(\beta\) can
be eliminated.

\Needspace{6\baselineskip}\textbf{Proposition 5.1.} For fixed \(n\geq2\) and \(p>1\), the
quantities \(c\) and \(k\) in (4.1) and (4.7), computed on a ball with
Robin parameter \(\beta>0\), satisfy

\[1\leq c\leq C_0(n,p)<\infty,
\qquad 0<k_0(n,p)\leq k\leq1.\tag{5.1}\]

Consequently the supremum of \(q_0\) over all radii and all positive
Robin parameters is finite.

\textbf{Proof.} Under \(x=Ry\), the dimensionless boundary parameter is
\(\alpha=\beta R^{p-1}\). The eigenvalue scales as \(R^{-p}\), the
normalized eigenfunction as \(R^{-n/p}\), and \(\tau\) as \(R^{n+p'}\).
It follows directly from (3.5), (4.1) and (4.7) that \(c\) and \(k\) are
unchanged by this simultaneous rescaling. We may therefore take \(R=1\)
and vary \(\beta\) in \((0,\infty)\).

Write \(V=|B_1|\). Integrating (3.1) gives

\[\tau(B_1)=V\left[b+\frac{n^{-1/(p-1)}}{n+p'}\right],
\qquad b=(n\beta)^{-1/(p-1)}.\tag{5.2}\]

Let \(\lambda_D\) be the first Dirichlet \(p\)-eigenvalue on \(B_1\).
The Rayleigh principle, using a Dirichlet eigenfunction and a constant
function respectively, gives

\[0<\lambda_B\leq\lambda_D,\qquad\lambda_B\leq n\beta.\tag{5.3}\]

Equation (3.2) implies

\[0\leq-\phi'(r)\leq\lambda_D^{1/(p-1)}\phi(0)r^{1/(p-1)}.\tag{5.4}\]

Choose \(r_*\in(0,1)\) depending only on \(n,p\) so that integration of
(5.4) yields \(\phi(r)\geq\phi(0)/2\) on \([0,r_*]\). Normalization then
bounds \(\phi(0)\) uniformly above, while monotonicity gives
\(\phi(0)\geq V^{-1/p}\). The normalized radial eigenfunctions are
therefore uniformly bounded and equicontinuous on \([0,1]\).

As \(\beta\downarrow0\), (5.3) and (3.2) show that the oscillation of
\(\phi\) tends to zero. Normalization implies uniform convergence to
\(V^{-1/p}\). The integrated eigenfunction equation is

\[\lambda_B\int_{B_1}\phi^{p-1}\,dx
=\beta|\partial B_1|\phi(1)^{p-1}.\tag{5.5}\]

Hence \(\lambda_B/(n\beta)\to1\). Inserting this limit and (5.2) into
the definitions of \(c\) and \(k\) gives

\[c\longrightarrow p,\qquad k\longrightarrow1
\qquad(\beta\downarrow0).\tag{5.6}\]

For the other endpoint, (5.4) and the Robin condition imply
\(\phi(1)=O(\beta^{-1/(p-1)})\). Given any sequence \(\beta\to\infty\),
take a subsequence along which \(\lambda_B\) converges to \(\lambda_*\)
and \(\phi\) converges uniformly to \(\phi_*\). The limit is
nonnegative, normalized, and vanishes at \(r=1\). Equation (3.2) passes
to the limit. Its right-hand side is bounded by a constant times
\(r^{1/(p-1)}\), so dominated convergence also gives convergence of the
gradient energies. Furthermore,

\[\beta\int_{\partial B_1}\phi^p\,d\mathcal H^{n-1}
=O(\beta^{-1/(p-1)})\longrightarrow0.\]

Therefore \(\phi_*\in W^{1,p}_0(B_1)\) and

\[\lambda_D\leq\int_{B_1}|\nabla\phi_*|^p\,dx
=\lambda_*\leq\lambda_D.\]

Thus \(\lambda_*=\lambda_D\) and \(\phi_*\) is the normalized
nonnegative first Dirichlet eigenfunction \(\phi_D\). Simplicity
identifies every subsequential limit, so the whole family converges.
Writing \(\tau_D=\int_{B_1}w_D\,dx\), we obtain

\[\begin{aligned}
c&\longrightarrow p\lambda_D^{1/(p-1)}\phi_D(0)^p\tau_D<\infty,\\
k&\longrightarrow\frac{\lambda_D\tau_D^{p-1}}{V^{p-1}}>0
\qquad(\beta\to\infty).
\end{aligned}\tag{5.7}\]

For completeness, the ball data are continuous at every
\(\beta_*\in(0,\infty)\). The same radial compactness and gradient
convergence apply along \(\beta\to\beta_*\). The limiting normalized
trial gives \(\lambda_*\geq\lambda_{\beta_*,p}(B_1)\). Conversely, the
fixed ground state at \(\beta_*\) is an admissible trial for nearby
parameters and gives the reverse upper limit. Simplicity again
identifies the uniform eigenfunction limit. Thus \(\lambda_B\) and
\(\phi(0)\), and hence \(c\) and \(k\), are continuous.

The endpoint limits and continuity prove (5.1), since \(c\geq1\) and
\(k\leq1\) were proved in Section 4. The denominator \(-\log(1-1/(2c))\)
in (4.7) is bounded below by the positive constant
\(-\log(1-1/(2C_0))\). Hence \(q_0\) has a finite supremum. \(\square\)

\textbf{Proof of Theorems 1.1 and 1.2.} Set

\[Q(n,p)=1+\sup_{\alpha>0}q_0(n,p,\alpha),\qquad
C(n,p)=2\sup_{\alpha>0}c(n,p,\alpha),\tag{5.8}\]

where the ball has radius one and boundary parameter \(\alpha\).
Proposition 5.1 makes these constants finite. Equation (4.5) proves
Theorem 1.2; Proposition 4.1 proves (1.5) and (1.6). If equality holds
in (1.5), (1.6) gives \(\theta=1\). Equation (4.4) then yields
\(\lambda(\Omega)\leq\lambda(B)\); Faber-Krahn and its rigidity
statement give that \(\Omega\) is a translate of \(B\). Such a translate
plainly attains equality. \(\square\)

\section{6. Consequences and remaining
questions}\label{consequences-and-remaining-questions}

Theorem 1.1 turns any available lower bound for the normalized torsional
deficit into the same lower bound for the product deficit. For example,
in the planar linear case the quantitative Robin Saint-Venant inequality
gives a positive constant \(a(V,\beta)\) such that

\[1-\mathcal R_q(\Omega)\geq a(V,\beta)\mathcal A(\Omega)^2
\qquad(q\geq Q(2,2)),\tag{6.1}\]

where
\(\mathcal A(\Omega)=\inf_{x\in\mathbb R^2}|\Omega\mathbin\triangle(B+x)|/V\)
is Fraenkel asymmetry. This is a direct consequence of (1.6) and the
known torsional remainder, as given for example in {[}ABCMP, Corollary
4.1{]}. The geometric estimate is a prior input; we do not claim that
\(a(V,\beta)\) can be chosen uniformly in \(\beta\) by that argument.

The proof of the product theorem uses no compactness of the class of
competing domains and no boundary perturbation of the ball. The
composition identity preserves the reference energy, and pointwise
torsion comparison controls the loss in the trial denominator. The
near-optimal estimate (4.4) and the global bound (4.6) are the two
distinct ingredients needed to reach every domain.

The same composition construction with zero boundary value is compatible
with the Dirichlet problem. The linear Dirichlet reverse inequality is
already a theorem of {[}BLNP{]}; it is not an additional open-problem
solution asserted here. In the Robin setting, the linear case \(n\geq3\)
remains beyond the pointwise input (1.8). The known \(L^1\) and \(L^p\)
torsion comparisons do not by themselves control the nonlinear numerator
\(\int H(w)\) in (4.2). Establishing a suitable substitute, or a
different trial construction, remains necessary to extend the present
proof.

The exponent (5.8) is sufficient and is not asserted to be optimal.
Determining the least admissible exponent, extending the range of \(p\),
and settling the forward Robin Kohler-Jobin question {[}BCS, Q1{]}
remain separate problems.

\sloppy\section{References}\label{references}

{[}ABCMP{]} V. Amato, R. Barbato, S. Cito, A. L. Masiello and G. Paoli,
\emph{A quantitative Talenti-type comparison result with Robin boundary
conditions}, preprint,
\href{https://arxiv.org/abs/2511.11316v1}{arXiv:2511.11316v1}.

{[}AGM{]} V. Amato, A. Gentile and A. L. Masiello, \emph{Comparison
results for solutions to p-Laplace equations with Robin boundary
conditions}, Annali di Matematica Pura ed Applicata \textbf{201} (2022),
1189-1212. \href{https://doi.org/10.1007/s10231-021-01153-y}{DOI:
10.1007/s10231-021-01153-y}.

{[}ANT{]} A. Alvino, C. Nitsch and C. Trombetti, \emph{A Talenti
comparison result for solutions to elliptic problems with Robin boundary
conditions}, Communications on Pure and Applied Mathematics \textbf{76}
(2023), 585-603. \href{https://doi.org/10.1002/cpa.22090}{DOI:
10.1002/cpa.22090}.

{[}BCS{]} G. Buttazzo, S. Cito and F. Solombrino, \emph{Relations
between principal eigenvalue and torsional rigidity with Robin boundary
conditions}, Milan Journal of Mathematics \textbf{94} (2026), 369-386.
\href{https://doi.org/10.1007/s00032-026-00438-2}{DOI:
10.1007/s00032-026-00438-2}.

{[}BD{]} D. Bucur and D. Daners, \emph{An alternative approach to the
Faber-Krahn inequality for Robin problems}, Calculus of Variations and
Partial Differential Equations \textbf{37} (2010), 75-86.
\href{https://doi.org/10.1007/s00526-009-0252-3}{DOI:
10.1007/s00526-009-0252-3}.

{[}BLNP{]} D. Bucur, J. Lamboley, M. Nahon and R. Prunier, \emph{Sharp
quantitative stability of the Dirichlet spectrum near the ball},
preprint, \href{https://arxiv.org/abs/2304.10916v2}{arXiv:2304.10916v2}.

\end{document}
